\documentclass{article}

\usepackage[polish]{babel}
\usepackage[a4paper,top=2cm,bottom=2cm,left=3cm,right=3cm,marginparwidth=1.75cm]{geometry}
\usepackage{amsmath}
\usepackage{graphicx}
\usepackage[colorlinks=true, allcolors=blue]{hyperref}
\usepackage[T1]{fontenc}

\title{Wyznaczanie aproksymacji układu lokalnego przez ukośne odwzorowanie Merkatora}

\begin{document}
\maketitle

\section{Koncepcja}
Celem omawianego narzędzia jest wyznaczenie definicji \href{https://proj.org/en/stable/operations/projections/omerc.html}{ukośnego odwzorowania Merkatora} 
aproksymującego lokalny układ współrzędnych kartezjańskich. Lokalne układy nadal często wykorzystywane są w przypadku obsługi geodezyjnej takich obiektów 
jak np. inwestycje budowlane lub place budowy. Układ lokalny ma wiele przewag w przypadku obsługi mniejszych powierzchniowo obszarów. Przez wyraźne zdefiniowanie 
jego początku (punktu przyłożenia) oraz osi możliwe jest między innymi łatwiejsze interpretowanie wartości współrzędnych w codziennych pracach.

Standardowo pracując z układem lokalnym konieczne jest określenie reguł matematycznych umożliwiających ścisłe przeliczanie pomiędzy współrzędnymi w układzie lokalnym, 
a innym globalnym układem współrzędnych. Rozwiązanie to choć najbardziej poprawne merytorycznie znacznie utrudnia pracę z danymi geodezyjnymi w środowisku GIS. 
W każdym wypadku kiedy potrzebny jest import danych do oprogramowania GIS konieczne jest przeliczenia współrzędnych lokalnych poza nim (i potencjalnie w kierunku odwrotnym 
- z układu globalnego do lokalnego). Z tego powodu alternatywnym rozwiązaniem jest zdefiniowanie poprawnego z punktu widzenia GIS układu odwzorowania będącego aproksymacją 
układu lokalnego, które dla odpowiednio małych obszarów nie będzie generowało żadnych różnic w wartościach współrzędnych.

\section{Ogólny opis procesu}
Generalny zarys procesu można zamknąć w dwóch głównych krokach - wyznaczeniu parametrów odwzorowania (punkt przyłożenia oraz azymut skręcenia), a następnie utworzeniu definicji układu.

Przewidujemy dwa warianty wyznaczania parametrów odwzorowania. W pierwszym z nich parametry wyznaczamy na podstawie dwóch punktów, gdzie jeden służy jako punkt przyłożenia, 
a drugi definiuje kierunek osi Y nowego układu. W drugim wariancie na podstawie większej grupy punktów w procesie wyrównania na podstawie wielu odpowiadających sobie par współrzędnych 
w układzie lokalnym i globalnym wyznaczane są parametry definiujące odwzorowanie. W wyniku obu tych wariantów finalnie otrzymujemy zestaw 3 parametrów określających definicję układu odwzorowania.

Znając wartości wszystkich wymaganych parametrów w następnym kroku definiujemy z wykorzystaniem PROJ/PyPROJ nowe odwzorowanie, 
które może być potem wykorzystywane do przeliczania współrzędnych, transformacji rastrów etc. w programach i bibliotekach do danych GIS.


\section{Kroki pośrednie}
\subsection{Definicja odwzorowania}
Żeby poprawnie zdefiniować odwzorowanie ukośne Merkatora konieczna jest znajomość następujących parametrów:
\begin{enumerate}
\item Punktu przyłożenia walca (2 wartości - szerokość i długość geograficzna),
\item Azymutu skręcenia walca (1 wartość kątowa)
\item Sztucznego przesunięcia do wartości współrzędnych XY (2 wartości)
\end{enumerate}
Należy nadmienić jednak, że musimy zdefiniować elipsoidę na której oparty jest układ. W naszym przypadku można wybierać pomiędzy
już istniejącymi definicjami WGS84 i GRS80 (w teorii minimalnie się różnią).

Znając podane parametry żeby zdefiniować odwzorowanie można skorzystać z jednej z dostępnych w PROJ metod definiowania układów.
Zgodnie z ich dokumentacją w naszej implemntacji wykorzystujemy standard OGC tzw. \href{https://proj.org/en/stable/faq.html#what-is-the-best-format-for-describing-coordinate-reference-systems}{CRS WKT-2}.

Sama definicja sprowadza się więc do wstawieniu w tekst z definicją układu odpowiednich wartości numerycznych, a następnie w razie potrzeby konwersja tej definicji do innych formatów opisu (np. PROJ4JS 
dla portali mapowych lub ESRI WKT dla oprogramowania ArcGIS). 

\subsection{Wyznaczanie parametrów odwzorowania z 2 punktów}
Znając dwa punkty o znanych współrzędnych w układzie lokalnym oraz układzie globalnym (docelowo we współrzędnych geodezyjnych na elipsoidzie) możemy wyznaczyć wszystkie z wymienionych wcześniej parametrów odwzorowania.

W pierwszej kolejności musimy ustalić rolę każdego z punktów: 
\begin{itemize}
    \item Punkt O - punkt który będzie stanowił punkt przyłożenia walca dla odwzorowania Merkatora. Jego położenie w układzie geodezyjnym zostanie podane jako punkt przyłożenia walca, 
    a wartości współrzędnych w układzie lokalnym posłużą jako wartości dla sztucznego przesunięcia XY.
    \item Punkt D - punkt który określi kierunek osi Y nowo definiowanego układu. 
\end{itemize}

Wartości 4 z pięciu potrzebnych parametrów w zasadzie wyznaczane są bezpośrednio z wartości współrzędnych punktu O. W celu wyznaczenia brakującego parametru (azymutu skręcenia) 
musimy wyznaczyć azymut z punktu O do punktu D. W tym celu trzeba rozwiązać odwrotny problem geodezyjny na elipsoidzie (wyznaczenie długości oraz azymutu linii geodezyjnej 
pomiędzy dwoma znanymi punktami). \href{https://proj.org/en/stable/faq.html#what-is-the-best-format-for-describing-coordinate-reference-systems}{Implementacja tego zadania} jest dostępna w PROJ/PyPROJ.

% \subsection{Wyznaczanie parametrów odwzorowania z N punktów}
% \begin{itemize}
% \item Wyznaczenie centroidu punktów w WGS,
% \item Wyznaczenie układu ENU dla centroidu.
% \item 
% \end{itemize}
\end{document}